\documentclass{article}
\usepackage{pathfinder-paper}

\title{When a Donor Wrist Cannot Locate a Rotating Held Object}

\begin{document}
\maketitle

\begin{abstract}
DynaMAC proposes the donor arm's forward-kinematic pose as a task frame when a held object is occluded during handover.
Aero Hand Open demonstrates rotation of a cube within a grasp using a tendon-driven hand. Reading these results together
gives a conditional limit: if the object-to-wrist transform changes, the wrist pose alone does not determine the
object's orientation. If two trials also have indistinguishable motor-encoder and action histories but require different
receiving approaches, no receiver restricted to those observations can guarantee the correct approach in both. The
papers provide a physical mechanism and the handover claim, but no matched-history object-pose measurements or DynaMAC
handover test on this hand. We specify the measurement needed to decide whether the limit has practical force.
% ledger: 6,7,9,10,11,13,14,17,18
\end{abstract}

\section{Introduction}

DynaMAC (paper P) treats the other arm's end-effector pose as a moving task frame for bimanual coordination. In its
handover discussion, it calls the held-object pose and donor-arm pose redundant for the receiver, and proposes the donor
arm's precise forward kinematics when object tracking suffers occlusion \cite{p}. % ledger: 6,9
Aero Hand Open (paper Q) provides a hand with sixteen joints, seven motors and only motor-encoder proprioception. Its
learned policy rotates a cube within the grasp in simulation and on hardware \cite{q}. % ledger: 6,7,9,10

The question here is whether P's wrist frame remains an adequate substitute for object pose if Q's hand rotates the
object while the wrist remains fixed. This concerns a particular claim about information at handover, rather than a
tested combination of the two systems. Their simulation and action interfaces differ, and neither paper reports this
handover. % ledger: 6,9,11,13

\section{What the source papers establish}

P models arm end-effector poses and uses the opposite arm as a candidate task parameter. Its account of handover
explicitly appeals to redundancy between donor-arm and held-object poses under occlusion \cite{p}. P disregards gripper
actions in its pose-level exposition for brevity; this says nothing about whether its implementation could command Q's
hand. % ledger: 6,9,10
Q reports in-hand cube rotation with hardware observations restricted to the seven motor encoders and the previous
action. The physical cube's pose trace is not logged, so Q establishes visible rotation, not a measured distribution of
object poses conditional on encoder and action histories \cite{q}. % ledger: 6,7,9,10

\section{Conditional limit of the wrist frame}

Let \(T_w\) denote donor wrist pose, \(T_o\) object pose and \(G\) the object pose relative to the wrist. By definition, % ledger: 11
\[
T_o=T_wG.
\]
For a receiving task, the wrist can stand in for object pose when \(G\) remains sufficiently stable for the required
approach. Q's in-hand rotation gives a concrete way for \(G\) to change while \(T_w\) stays fixed. It does not quantify
how uncertain \(G\) is at a given motor history. Thus the result is a limit on P's stated substitution, not evidence of
a DynaMAC failure. % ledger: 6,9,10,11,13

\begin{proposition}[Matched-observation limit]
Suppose two occluded handover trials have indistinguishable donor wrist poses, motor-encoder histories and action
histories at the receiver's decision, but their hidden object orientations require different receiving approaches. A
receiver using only those observations cannot guarantee an orientation-specific successful approach in both trials.
% ledger: 7,14,17
\end{proposition}

\begin{proof}
The receiver receives the same information in both trials. A deterministic selection is therefore the same in both; a
randomised selection has the same distribution in both. Neither can guarantee two different required approaches from
that information. % ledger: 7,14,17
\end{proof}

This proposition is conditional on the existence of such matched trials and an orientation-sensitive receiving task. Q's
seven encoders and its rotation demonstration do not establish that naturally occurring trials have indistinguishable
full histories but different object poses. Controlled reseating under occlusion is one way to construct candidate pairs;
its frequency in ordinary manipulation is unknown. If one safe grasp works for both orientations, the proposed approach
distinction disappears. % ledger: 7,9,10,14,17

\section{Related work and interpretation}

The main step has a precedent in pose-aware handover: Vezzani and colleagues estimate in-hand object pose from vision
and tactile information and use it to choose the receiving hand's pose \cite{vezzani2017}. Liang and colleagues track a
moving in-hand object's pose under visual occlusion using contact feedback and simulation \cite{liang2020}. Object-pose
estimation and pose-aware receiving are therefore already published. The contribution here is the conditional,
pair-specific reading of P's donor-wrist substitution against Q's in-hand rotation, together with a discriminating test.
We do not claim that this logical limit or a remedy is novel in general. % ledger: 11,12,13,15,16,18

P's kinematic-link detector presents a separate question. Its geometric mean standard deviation describes dispersion of
a learned object-frame policy stream across phase-aligned demonstrations, not online variation in \(G\) during one
rollout. In-hand motion or compliance alone therefore does not show that the detector would retain or mask the object
frame incorrectly. Testing detector decisions against controlled contact or motor-perturbation labels remains distinct
from the wrist-frame claim \cite{p}. % ledger: 3,5,8,18

An uncertainty-aware receiver faced with exactly matched histories under full occlusion could represent the ambiguity,
seek a new view, or use a grasp robust to both orientations. It could not recover which hidden orientation occurred
without new information. These are possible responses to the limit, not methods evaluated here. % ledger: 14,17

\section{A decisive test}

Hold the donor wrist fixed and externally measure \(G\) for a marked object while varying its in-hand orientation or
reseated contact state. First ask whether materially different object poses occur within matched bins of wrist pose,
motor readings and action history. Then, only if they do, use an orientation-sensitive receiving target and check that
no single orientation-invariant grasp succeeds across the matched conditions. Compare a receiver using only the donor
wrist frame with visible-object and last-seen-object controls; a rigid-grasp or no-reorientation condition would help
isolate the role of \(G\). Occlude the object at receiver inference while retaining external pose measurement for
analysis. This is a proposed experiment, not an outcome reported by either paper. % ledger: 7,11,14,17,18

\section{Limitations and open questions}

No physical trace of Q's cube pose, conditional spread of \(G\), detector error rate, or handover success with DynaMAC
on Q's hardware is available. Constructing matched encoder and action histories may require external reseating, which
might not represent routine use. The papers' distinct simulators and action spaces make a direct benchmark swap
substantial work. The practical question is whether measured ambiguity is large enough to change a receiving decision,
and whether obtaining another observation improves that decision. Neither answer follows from the current material.
% ledger: 5,9,10,13,14,17,18

\bibliographystyle{plainurl}
\bibliography{references}
\end{document}
